% 出典: 2025ap_exam.pdf, 日本語版 PDF pp. 2--3.
\documentclass[11pt,a4paper]{article}
\usepackage[margin=24mm]{geometry}
\usepackage{fontspec}
\usepackage{xeCJK}
\setCJKmainfont{HaranoAjiMincho-Regular.otf}
\setCJKsansfont{HaranoAjiGothic-Medium.otf}
\usepackage{amsmath,amssymb}
\usepackage{tikz}
\usetikzlibrary{arrows.meta}
\setlength{\parindent}{1em}
\setlength{\parskip}{0.35em}
\title{物理工学専攻 院試解答集\\2025年 物理学 第1問}
\author{政岡凜太郎}
\date{}
\begin{document}
\maketitle

\section*{[1] 糸で落下する剛体}
図1の内側の円盤は厚さを無視するので、剛体の質量は外側の2枚に分布する。一方、糸は半径 $r_2$ の内側の円盤に接している。この二つの半径の役割を区別する。

\begin{center}
\begin{tikzpicture}[>=Stealth,scale=0.92]
  \draw[thick] (0,0) circle (1.65);
  \draw[dashed] (0,0) circle (0.78);
  \fill (0,0) circle (1.3pt) node[below left] {$C$};
  \draw[very thick] (0.78,0)--(0.78,2.1);
  \draw[->] (0,0)--(1.18,1.16) node[midway,above left] {$r_1$};
  \draw[->] (0,0)--(0.78,0) node[midway,below] {$r_2$};
  \draw[->,thick] (2.25,1.25)--(2.25,0.28) node[midway,right] {$a$};
  \node[above] at (0.78,2.1) {糸の固定端};
  \node[below] at (0,-1.9) {外円盤の半径と糸のトルク腕};
\end{tikzpicture}
\end{center}

\subsection*{[1.1] 中心軸まわりの慣性モーメント}
外側の2枚は同じ厚さ $d_1$ と同じ密度を持つため、1枚の質量は $M/2$ である。1枚について半径 $s$ から $s+\mathrm ds$ の環の質量は、面密度を $(M/2)/(\pi r_1^2)$ として
\[
\mathrm dm=\frac{M/2}{\pi r_1^2}\,2\pi s\,\mathrm ds
=\frac{M}{r_1^2}s\,\mathrm ds.
\]
よって1枚の中心軸まわりの慣性モーメントは $\int_0^{r_1}s^2\,\mathrm dm=Mr_1^2/4$。2枚を足すと
\[
\boxed{I_C=\frac12Mr_1^2}.
\]
厚さを無視する内側の円盤には質量を割り当てない。$d_1$ は面密度を通じて $M$ に含まれている。

\subsection*{[1.2] 並進と回転の運動方程式}
並進の正方向を鉛直下向き、回転の正方向を糸がほどける向きに選ぶ。重力と糸の張力による並進の式は
\[
\boxed{Ma=Mg-T}.
\]
糸の張力は中心軸から距離 $r_2$ で作用するので、角速度のこの正方向の成分を $\omega$ とすれば
\[
\boxed{I_C\dot\omega=Tr_2}.
\]
糸が滑らないため、糸のほどけた長さは重心の下降距離に等しい。これを時間で二度微分すると $\boxed{a=r_2\dot\omega}$ となる。

\subsection*{[1.3] 重心の加速度}
回転の式と滑りなしの条件から $T=I_Ca/r_2^2$。これを並進の式に代入し、[1.1] の $I_C$ を使うと
\[
\left(M+\frac{Mr_1^2}{2r_2^2}\right)a=Mg,
\qquad
\boxed{a=\frac{2r_2^2}{r_1^2+2r_2^2}\,g}.
\]
$r_1,r_2,g>0$ なら $0<a<g$ であり、$T=M(g-a)>0$ とも整合する。

\section*{[2] 磁場中の帯電球殻}
位置 $\boldsymbol r$ は球殻の重心から表面上の点へ向かう。電荷は球殻とともに回るので、各表面素片の電荷に働く磁気力を求め、そのトルクを球面全体で積分する。外部磁場 $\boldsymbol B=(0,0,B)$ だけを用いる。

\subsection*{[2.1] 表面上の速度}
角速度 $\boldsymbol\omega$ で剛体回転する点の速度は、中心が静止しているため
\[
\boxed{\boldsymbol v=\boldsymbol\omega\times\boldsymbol r}.
\]

\subsection*{[2.2] 微小面積に働く力}
面積 $\mathrm dS$ の電荷は $\sigma\,\mathrm dS$。磁場によるローレンツ力は、問題の三重積公式に $\boldsymbol a=\boldsymbol\omega$, $\boldsymbol b=\boldsymbol r$, $\boldsymbol c=\boldsymbol B$ を代入して
\[
\boxed{\mathrm d\boldsymbol F
=\sigma\,\mathrm dS\,(\boldsymbol\omega\times\boldsymbol r)\times\boldsymbol B
=\sigma\,\mathrm dS\left[(\boldsymbol\omega\cdot\boldsymbol B)\boldsymbol r
-(\boldsymbol r\cdot\boldsymbol B)\boldsymbol\omega\right]}.
\]
$\sigma$ は符号付きの面電荷密度なので、負電荷では力の向きが反転する。

\subsection*{[2.3] 微小トルク}
原点まわりのトルクは $\boldsymbol r\times\mathrm d\boldsymbol F$。力のうち $\boldsymbol r$ と平行な項は寄与しない。したがって
\[
\boxed{\mathrm d\boldsymbol N
=-\sigma\,\mathrm dS\,(\boldsymbol r\cdot\boldsymbol B)
   (\boldsymbol r\times\boldsymbol\omega)
=\sigma\,\mathrm dS\,(\boldsymbol r\cdot\boldsymbol B)
   (\boldsymbol\omega\times\boldsymbol r)}.
\]

\subsection*{[2.4] 球面全体からのトルク}
球面の回転・反射対称性［問題の式 (1)］により、$\int_{\Omega_R}x^2\,\mathrm dS$, $\int_{\Omega_R}y^2\,\mathrm dS$, $\int_{\Omega_R}z^2\,\mathrm dS$ は等しく、異なる座標の積の積分はゼロである。$x^2+y^2+z^2=R^2$ と球面積 $4\pi R^2$ から
\[
\int_{\Omega_R}x^2\,\mathrm dS
=\int_{\Omega_R}y^2\,\mathrm dS
=\int_{\Omega_R}z^2\,\mathrm dS
=\frac{4\pi R^4}{3}.
\]
よって $\int_{\Omega_R}(\boldsymbol r\cdot\boldsymbol B)\boldsymbol r\,\mathrm dS=(4\pi R^4/3)\boldsymbol B$。$\boldsymbol\omega$ は球面上の位置に依らないので [2.3] を積分して
\[
\boldsymbol N=\frac{4\pi\sigma R^4}{3}\,\boldsymbol\omega\times\boldsymbol B.
\]
$\boldsymbol B=(0,0,B)$ を代入した答えは
\[
\boxed{N_x=\frac{4\pi\sigma BR^4}{3}\omega_y,\qquad
N_y=-\frac{4\pi\sigma BR^4}{3}\omega_x,\qquad N_z=0}.
\]

\subsection*{[2.5] 各軸の慣性モーメント}
面質量密度 $\rho$ について $x$ 軸からの距離の二乗は $y^2+z^2$。したがって
\[
I_x=\rho\int_{\Omega_R}(y^2+z^2)\,\mathrm dS
=\frac{8\pi\rho R^4}{3}.
\]
同じ球対称性を $y,z$ 軸にも用いると
\[
\boxed{I_x=I_y=I_z=I=\frac{8\pi\rho R^4}{3}}.
\]

\subsection*{[2.6] 角速度と歳差の向き}
球殻の慣性モーメントはどの軸でも $I$ なので、角運動量は $I\boldsymbol\omega$ であり、回転運動方程式は $I\dot{\boldsymbol\omega}=\boldsymbol N$。[2.4] と [2.5] を代入して
\[
\dot\omega_x=\frac{\sigma B}{2\rho}\omega_y,\qquad
\dot\omega_y=-\frac{\sigma B}{2\rho}\omega_x,\qquad
\dot\omega_z=0.
\]
初期条件は $\omega_x(0)=\omega_0\sin\theta$, $\omega_y(0)=0$, $\omega_z(0)=\omega_0\cos\theta$。最初の二式から $\ddot\omega_x=-(\sigma B/2\rho)^2\omega_x$ となり、初期値を満たす解は
\[
\boxed{\begin{aligned}
\omega_x(t)&=\omega_0\sin\theta\cos\!\left(\frac{\sigma B}{2\rho}t\right),\\
\omega_y(t)&=-\omega_0\sin\theta\sin\!\left(\frac{\sigma B}{2\rho}t\right),\\
\omega_z(t)&=\omega_0\cos\theta.
\end{aligned}}
\]
横成分を$\boldsymbol\omega_\perp=(\omega_x,\omega_y,0)$と書く。その大きさは $\omega_0\sin\theta$ のまま $xy$ 平面を回り、$\omega_z$ も一定である。したがって回転軸は $z$ 軸となす角 $\theta$ を保って歳差する。

\begin{center}
\begin{tikzpicture}[>=Stealth,scale=0.9]
  \draw[->] (-0.4,0)--(2.6,0) node[right] {$+x$};
  \draw[->] (0,-0.35)--(0,2.2) node[above] {$+y$};
  \draw[dashed] (0,0) circle (1.65);
  \draw[->,very thick] (0,0)--(1.65,0) node[below right] {$\boldsymbol\omega_{\perp}(0)$};
  \draw[->,thick] (1.7,0.16) arc[start angle=5,end angle=65,radius=1.7];
  \node[anchor=west] at (2.0,1.25) {$\sigma<0,\ B>0$};
  \node[below] at (0,-2.0) {$+z$ 側から見た横成分};
\end{tikzpicture}
\end{center}

$\sigma<0,B>0$ なら、初期時刻の $\dot\omega_y(0)=-\sigma B\omega_0\sin\theta/(2\rho)>0$。すなわち横成分は $+x$ から $+y$ に向かうので、$+z$ 側から見た歳差の向きは\textbf{反時計回り}である。
\end{document}
